% Source: https://artofproblemsolving.com/community/c3h2519584_two_three_digit_numbers_with_difference_slovenia_national_mo_1999_1st_grade_p1
% Stable identity: aops-topic-2519584-q01
Let the larger number be $X=\overline{pqr}$. Its units digit $r$ is the hundreds digit of the smaller number, and its hundreds digit $p$ is the units digit of the smaller one, so the smaller number is $Y=\overline{rsp}$ for some digit $s$, with $p,r\ge1$. Then
\[X-Y=(100p+10q+r)-(100r+10s+p)=99(p-r)+10(q-s)=297.\]
Since $|q-s|\le9$, we have $99(p-r)=297-10(q-s)\in[207,387]$, and the only multiple of $99$ in this interval is $297$. So $p-r=3$ and $q=s$.

The digit sum of $Y$ is $r+s+p=r+q+(r+3)=23$, so $2r+q=20$. Since $p=r+3\le9$, $r\le6$; since $q\le9$, $2r\ge11$, so $r\ge6$. Hence $r=6$, $q=8$, $p=9$, and the numbers are
\[\boxed{986\text{ and }689}.\]
Check: $986-689=297$ and $6+8+9=23$.
